\subsection{Методы, не требующие производной}
\label{subsec:direct-search-1d}

Все методы ниже устроены по одной схеме: вычислить $f$ в нескольких пробных
точках, применить лемму~\ref{lem:unimodal} и выбросить часть отрезка. Различие
только в том, \emph{где} ставить пробные точки, — и именно оно определяет цену
метода в числе обращений к $f$.

\begin{idea}[title={О форме записи алгоритмов}]
В тетради алгоритмы записаны схематично: для большинства методов выписано
только правило сокращения отрезка, без явных формул инициализации и критерия
остановки. Ниже каждый метод оформлен как полная пошаговая процедура; места,
где запись в тетради расходится с логикой леммы, отмечены отдельно.
\end{idea}

\subsubsection*{Метод равномерного поиска}

Отрезок разбивается равномерной сеткой, число разбиений $n$ задано заранее.

\begin{algo}[title={Метод равномерного поиска}]
\textbf{Вход:} $[a,b]$, число разбиений $n$, требуемая точность $\delta$.

\begin{steps}
  \item Положить $h = \dfrac{b-a}{n}$ и $x_i = a + h\,i$, $i = \overline{0,n}$.
  \item Вычислить $f(x_i)$ для всех $i$ и найти номер $j$, на котором
        достигается минимум по сетке:
        \[
          f(x_j) = \min_{i} f(x_i) .
        \]
  \item Тогда
        \[
          \begin{aligned}
            f(x_{j-1}) &> f(x_j) &&\Longrightarrow&& x_* \in [x_{j-1},\, b] , \\
            f(x_{j+1}) &\ge f(x_j) &&\Longrightarrow&& x_* \in [a,\, x_{j+1}] ,
          \end{aligned}
        \]
        и, пересекая, получаем
        \[
          x_* \in [x_{j-1},\, x_{j+1}] .
        \]
  \item Если $2h \le \delta$ — остановиться. Иначе положить
        $[a,b] := [x_{j-1},\, x_{j+1}]$ и вернуться к шагу~1.
\end{steps}
\end{algo}

\inputfig{figures/fig-uniform-grid}

\begin{remark}
Если минимум по сетке достигается на краю ($j = 0$ или $j = n$), надо считать
$x_{-1} := a$ и $x_{n+1} := b$. Для унимодальной $f$ такой номер $j$
единственный.
\end{remark}

\begin{idea}[title={Какое $n$ брать — оценка (на лекции осталось вопросом)}]
За один проход интервал сокращается в $n/2$ раз, а стоит это примерно $n-1$
новых вычислений (значения на концах уже известны с прошлого прохода). Значит
«коэффициент сокращения на одно вычисление» равен
\[
  q(n) = \left(\frac{2}{n}\right)^{1/(n-1)} ,
\]
и его надо минимизировать:
\[
  q(3) \approx 0{,}816, \quad
  \mathbf{q(4) \approx 0{,}794}, \quad
  q(5) \approx 0{,}795, \quad
  q(6) \approx 0{,}803 .
\]
Оптимум — $n = 4$. Для сравнения: золотое сечение даёт $0{,}618$ на одно
вычисление, то есть равномерный поиск проигрывает ему при любом $n$. Выкладка
моя, на лекции её не было.
\end{idea}

\subsubsection*{Метод локализации экстремума}

Отличается от равномерного поиска тем, что не нужно вычислять все значения на
сетке: идём слева направо и останавливаемся на первом «провале».

\begin{algo}[title={Метод локализации экстремума}]
\textbf{Вход:} $[a,b]$, число разбиений $n$, точность $\delta$.

\begin{steps}
  \item Положить $h = \dfrac{b-a}{n}$.
  \item Последовательно для $i = 1, 2, \dots, n-1$ проверять пару условий
        \[
          \begin{cases}
            f\big(a + (i-1)h\big) > f\big(a + ih\big), \\[2pt]
            f\big(a + ih\big) \le f\big(a + (i+1)h\big),
          \end{cases}
        \]
        и остановиться на первом номере $i$, для которого они выполнены.
  \item Тогда
        \[
          x_* \in \big[\, a + (i-1)h, \ a + (i+1)h \,\big] .
        \]
  \item Если $2h \le \delta$ — остановиться. Иначе повторить процесс для
        отрезка $\big[a + (i-1)h,\ a + (i+1)h\big]$.
\end{steps}
\end{algo}

\begin{idea}
Для унимодальной функции найденный здесь номер $i$ — тот же самый, что и номер
$j$ в равномерном поиске: первый «провал» и есть минимум по сетке. Выигрыш в
том, что в среднем значений вычисляется меньше, чем $n+1$: до минимума функция
убывает, и обход прекращается сразу после него.
\end{idea}

\subsubsection*{Метод золотого сечения}

Здесь пробные точки ставятся симметрично и так, чтобы \textbf{одна из них
переиспользовалась} на следующем шаге.

На $k$-м шаге на отрезке $[a_k, b_k]$ берутся
\[
  \lambda_k = a_k + \alpha\,(b_k - a_k), \qquad
  \mu_k = a_k + b_k - \lambda_k = b_k - \alpha\,(b_k - a_k),
  \qquad \alpha \in \left(0,\tfrac12\right) ,
\]
так что выполняется симметрия
\[
  b_k - \lambda_k = \mu_k - a_k .
\]

\inputfig{figures/fig-golden-section}

Правило сокращения (лемма~\ref{lem:unimodal}):
\[
  f(\lambda_k) \ge f(\mu_k) \ \Longrightarrow\ [\lambda_k,\, b_k] ,
  \qquad\qquad (1)
\]
\[
  f(\lambda_k) < f(\mu_k) \ \Longrightarrow\ [a_k,\, \mu_k] .
  \qquad\qquad (2)
\]

\paragraph{Подбор $\alpha$.} Необходимо сократить число обращений к $f$, то
есть добиться
\[
  \lambda_{k+1} = \mu_k \quad\text{в случае (1)},
  \qquad
  \mu_{k+1} = \lambda_k \quad\text{в случае (2)} .
\]
Обозначим $L_k = b_k - a_k$. В случае~(1) имеем $a_{k+1} = \lambda_k$,
$b_{k+1} = b_k$, значит $L_{k+1} = (1-\alpha)L_k$ и
\[
  \lambda_{k+1} = a_{k+1} + \alpha L_{k+1}
  = \lambda_k + \alpha(1-\alpha)L_k
  = a_k + \alpha L_k + \alpha(1-\alpha)L_k .
\]
Приравнивая это к $\mu_k = a_k + (1-\alpha)L_k$, получаем
\[
  \alpha + \alpha(1-\alpha) = 1 - \alpha
  \quad\Longleftrightarrow\quad
  2\alpha - \alpha^2 = 1 - \alpha
  \quad\Longleftrightarrow\quad
  \alpha^2 - 3\alpha + 1 = 0 ,
\]
откуда
\[
  \alpha_{1,2} = \frac{3 \pm \sqrt5}{2},
  \qquad
  \boxed{\ \alpha = \frac{3 - \sqrt5}{2} \approx 0{,}382\ }
\]
(корень со знаком «$+$» больше единицы и не подходит).

\begin{idea}
Случай~(2) даёт то же самое уравнение: там $\mu_{k+1} = a_k + (1-\alpha)^2 L_k$
и условие $\mu_{k+1} = \lambda_k$ означает $(1-\alpha)^2 = \alpha$, то есть
снова $\alpha^2 - 3\alpha + 1 = 0$. Так что одним подбором $\alpha$
обслуживаются оба исхода.

\medskip
Коэффициент сокращения равен
\[
  1 - \alpha = \frac{\sqrt5 - 1}{2} \approx 0{,}618 ,
\]
и это ровно золотое сечение — отсюда название. На каждой итерации, кроме
первой, тратится \textbf{одно} новое вычисление функции.
\end{idea}

\begin{algo}[title={Метод золотого сечения}]
\textbf{Вход:} $[a_1,b_1] = [a,b]$, точность $\delta$,
$\alpha = \tfrac{3-\sqrt5}{2}$.

\begin{steps}
  \item Положить
        $\lambda_1 = a_1 + \alpha(b_1-a_1)$,
        $\mu_1 = b_1 - \alpha(b_1-a_1)$;
        вычислить $f(\lambda_1)$ и $f(\mu_1)$.
  \item Если $b_k - a_k \le \delta$ — остановиться, принять
        $x_* \approx \tfrac{a_k+b_k}{2}$.
  \item Если $f(\lambda_k) \ge f(\mu_k)$, положить
        \[
          a_{k+1} = \lambda_k, \quad b_{k+1} = b_k, \quad
          \lambda_{k+1} = \mu_k \ \text{(значение известно)},
        \]
        и вычислить одно новое значение
        $f(\mu_{k+1})$, где $\mu_{k+1} = b_{k+1} - \alpha(b_{k+1}-a_{k+1})$.
  \item Иначе положить
        \[
          a_{k+1} = a_k, \quad b_{k+1} = \mu_k, \quad
          \mu_{k+1} = \lambda_k \ \text{(значение известно)},
        \]
        и вычислить одно новое значение
        $f(\lambda_{k+1})$, где $\lambda_{k+1} = a_{k+1} + \alpha(b_{k+1}-a_{k+1})$.
  \item Увеличить $k$ и вернуться к шагу~2.
\end{steps}
\end{algo}

\subsubsection*{Метод дихотомии}

Пробные точки $x_1 < x_2$ берутся вплотную к середине отрезка. Правило
сокращения то же:
\[
  f(x_1) \ge f(x_2) \ \Longrightarrow\ [x_1,\, b] ,
  \qquad
  f(x_1) < f(x_2) \ \Longrightarrow\ [a,\, x_2] .
\]

\paragraph{Почему именно середина.} Длина оставшегося отрезка равна либо
$b - x_1$, либо $x_2 - a$; заранее неизвестно, какой исход реализуется,
поэтому точки выбираются из условия наилучшего худшего случая:
\[
  \min_{x_1,\, x_2} \ \max\big\{\, b - x_1,\ x_2 - a \,\big\} .
\]
Первая величина убывает по $x_1$, вторая растёт по $x_2$, поэтому минимум
максимума достигается при
\[
  x_1 = x_2 = \frac{a+b}{2} .
\]

\inputfig{figures/fig-dichotomy-minmax}

Совпадающие точки бесполезны, поэтому их разносят на малую величину
$\eps > 0$:
\[
  x_1 := \frac{a+b}{2} - \eps , \qquad
  x_2 := \frac{a+b}{2} + \eps .
\]

\begin{algo}[title={Метод дихотомии}]
\textbf{Вход:} $[a_1,b_1] = [a,b]$, точность $\delta$, параметр разведения
$\eps > 0$.

\begin{steps}
  \item Если $b_k - a_k \le \delta$ — остановиться.
  \item Положить
        $x_1 = \dfrac{a_k+b_k}{2} - \eps$, \ \
        $x_2 = \dfrac{a_k+b_k}{2} + \eps$;
        вычислить $f(x_1)$ и $f(x_2)$.
  \item Если $f(x_1) \ge f(x_2)$, положить $[a_{k+1},b_{k+1}] = [x_1,\, b_k]$;
        иначе $[a_{k+1},b_{k+1}] = [a_k,\, x_2]$.
  \item Увеличить $k$ и вернуться к шагу~1.
\end{steps}
\end{algo}

\begin{remark}[выбор $\eps$]
$\eps$ берётся настолько малым, насколько позволяет различимость значений
$f\!\left(\tfrac{a+b}{2} - \eps\right)$ и
$f\!\left(\tfrac{a+b}{2} + \eps\right)$: если они сольются в пределах точности
вычислений, сравнение потеряет смысл. Практический ориентир из тетради — доли
процента от рассматриваемых величин.
\end{remark}

\begin{idea}
Длина отрезка за итерацию становится равной $\tfrac{L_k}{2} + \eps$, то есть
почти вдвое меньше, но ценой \textbf{двух} вычислений функции. На одно
вычисление это даёт $\sqrt{1/2} \approx 0{,}707$ — хуже, чем $0{,}618$ у
золотого сечения.
\end{idea}

\subsubsection*{Метод Фибоначчи}

Та же идея симметричных точек с переиспользованием, но коэффициенты берутся не
постоянные, а из последовательности Фибоначчи
\[
  F_0 = F_1 = 1, \qquad F_{i+1} = F_i + F_{i-1} ,
\]
причём общее число шагов $n$ фиксируется \emph{заранее}.

На $k$-м шаге на отрезке $[a_k, b_k]$ берутся
\[
  \lambda_k := a_k + \frac{F_{n-k-1}}{F_{n-k+1}}\,(b_k - a_k) ,
  \qquad
  \mu_k := a_k + \frac{F_{n-k}}{F_{n-k+1}}\,(b_k - a_k) ,
\]
и $\lambda_k < \mu_k$, поскольку $F_{n-k-1} < F_{n-k}$. Правило сокращения
прежнее:
\[
  f(\lambda_k) \ge f(\mu_k) \ \Longrightarrow\ [\lambda_k,\, b_k] ,
  \qquad
  f(\lambda_k) < f(\mu_k) \ \Longrightarrow\ [a_k,\, \mu_k] .
\]

\begin{idea}[title={Проверка переиспользования — ответ на «???» в тетради}]
В тетради записано «$\mu_{k+1} = \lambda_k$ \textbf{???}». Проверяется прямой
подстановкой. Пусть $L_k = b_k - a_k$.

\medskip
\textbf{Случай (1):} $a_{k+1} = \lambda_k$, $b_{k+1} = b_k$, и, пользуясь
$F_{n-k+1} = F_{n-k} + F_{n-k-1}$,
\[
  L_{k+1} = L_k\left(1 - \frac{F_{n-k-1}}{F_{n-k+1}}\right)
          = L_k\,\frac{F_{n-k}}{F_{n-k+1}} .
\]
Тогда
\[
  \lambda_{k+1} = \lambda_k + \frac{F_{n-k-2}}{F_{n-k}}\,L_{k+1}
                = \lambda_k + L_k\,\frac{F_{n-k-2}}{F_{n-k+1}}
                = \mu_k ,
\]
поскольку $\mu_k - \lambda_k = L_k\,\dfrac{F_{n-k} - F_{n-k-1}}{F_{n-k+1}}
= L_k\,\dfrac{F_{n-k-2}}{F_{n-k+1}}$.

\medskip
\textbf{Случай (2):} $b_{k+1} = \mu_k$, $L_{k+1} = L_k\,\dfrac{F_{n-k}}{F_{n-k+1}}$,
и точно так же $\mu_{k+1} = \lambda_k$. То есть переиспользуется ровно одна
точка, как и в золотом сечении.
\end{idea}

\begin{algo}[title={Метод Фибоначчи}]
\textbf{Вход:} $[a_1,b_1] = [a,b]$, число шагов $n$, числа $F_0,\dots,F_n$.

\begin{steps}
  \item Положить $k = 1$, вычислить $f(\lambda_1)$ и $f(\mu_1)$ по формулам
        выше.
  \item Если $f(\lambda_k) \ge f(\mu_k)$, положить
        $[a_{k+1},b_{k+1}] = [\lambda_k,\, b_k]$ и $\lambda_{k+1} = \mu_k$
        (значение известно); вычислить $f(\mu_{k+1})$.
  \item Иначе положить $[a_{k+1},b_{k+1}] = [a_k,\, \mu_k]$ и
        $\mu_{k+1} = \lambda_k$ (значение известно); вычислить
        $f(\lambda_{k+1})$.
  \item Повторять, пока $k < n-1$.
\end{steps}
\end{algo}

\begin{remark}
На последнем шаге $k = n-1$ обе точки сливаются:
\[
  \lambda_{n-1} = \mu_{n-1} = \frac{a_{n-1} + b_{n-1}}{2} ,
\]
и сравнивать становится нечего. Поэтому одну из точек сдвигают на малое $\eps$
— ровно как в методе дихотомии. В тетради этот момент записан как «далее
возврат на предыдущий шаг».
\end{remark}

\begin{idea}
Итоговая длина интервала неопределённости равна
\[
  \frac{b-a}{F_n} ,
\]
и это \textbf{наилучший возможный} результат при заданном числе вычислений
функции — в этом смысле метод Фибоначчи оптимален. Золотое сечение — его
предельный вариант: отношение $F_{n-1}/F_n$ стремится к $0{,}618$, поэтому при
больших $n$ оба метода работают практически одинаково, но золотому сечению не
нужно знать $n$ заранее.
\end{idea}

\subsubsection*{Метод бисекции}

На отрезке $[a_k, b_k]$ берутся три точки, делящие его на четыре равные части:
\[
  x_1 = a_k + \frac{L_k}{4}, \qquad
  x_2 = \frac{a_k + b_k}{2}, \qquad
  x_3 = a_k + \frac{3L_k}{4},
  \qquad L_k = b_k - a_k .
\]

\begin{algo}[title={Метод бисекции}]
\textbf{Вход:} $[a_1,b_1] = [a,b]$, точность $\delta$.

\begin{steps}
  \item Если $b_k - a_k \le \delta$ — остановиться.
  \item Вычислить $f(x_1)$, $f(x_2)$, $f(x_3)$ в точках выше (значение в
        середине $x_2$ обычно уже известно с прошлой итерации).
  \item Если $f(x_1) < f(x_2)$, то $x_* \in [a_k,\, x_2]$ — взять новым
        отрезком $[a_k, x_2]$.
  \item Иначе $x_* \in [x_1,\, b_k]$, и сравнить $f(x_2)$ с $f(x_3)$:
        \begin{itemize}[leftmargin=1.4em, topsep=2pt]
          \item если $f(x_2) > f(x_3)$, то $x_* \in [x_2,\, b_k]$;
          \item если $f(x_2) \le f(x_3)$, то $x_* \in [x_1,\, x_3]$.
        \end{itemize}
  \item Увеличить $k$ и вернуться к шагу~1.
\end{steps}
\end{algo}

\begin{idea}
В каждом из трёх исходов новый отрезок вдвое короче прежнего, а его серединой
оказывается одна из уже вычисленных точек ($x_1$, $x_3$ или $x_2$
соответственно). Поэтому на итерацию приходится два новых вычисления функции —
столько же, сколько в дихотомии, и с тем же коэффициентом $0{,}707$ на
вычисление.
\end{idea}

\subsubsection*{Сводка}

\begin{center}
\begin{tabular}{lcc}
\toprule
Метод & сокращение за итерацию & на одно вычисление $f$ \\
\midrule
Равномерный поиск ($n=4$) & $1/2$ за $3$ вычисления & $\approx 0{,}794$ \\
Дихотомия                 & $\approx 1/2$ за $2$    & $\approx 0{,}707$ \\
Бисекция                  & $1/2$ за $2$            & $\approx 0{,}707$ \\
Золотое сечение           & $0{,}618$ за $1$        & $\approx 0{,}618$ \\
Фибоначчи                 & $(b-a)/F_n$ за $n$      & оптимум \\
\bottomrule
\end{tabular}
\end{center}

\begin{idea}[title={О сводной таблице}]
Таблица собрана мной по разобранным оценкам — на лекции её не было. Числа в
последнем столбце — коэффициент, во сколько раз сокращается интервал
неопределённости в пересчёте на одно обращение к $f$: чем меньше, тем лучше.
\end{idea}

\subsubsection*{Не разобрано}

\begin{idea}[title={Не разобрано на лекции}]
В конце тетради перечислены ещё два метода, оба со знаками вопроса, то есть на
лекции они не разбирались:
\begin{itemize}[leftmargin=1.6em, topsep=2pt]
  \item \textbf{метод пробных точек};
  \item \textbf{метод аппроксимации параболой}.
\end{itemize}
\end{idea}

\begin{idea}
Для ориентира: метод аппроксимации параболой — это единственный из списка,
использующий не только сравнения, но и сами значения. По трём точкам строится
квадратичная интерполяция, и за новое приближение берётся вершина параболы.
На гладкой функции вблизи минимума он сходится заметно быстрее любого метода
деления отрезка, но гарантий, как у перечисленных выше, не даёт.
\end{idea}
